\documentclass[10pt,a4paper]{amsart}
\usepackage{amsmath,amssymb,amsthm,graphicx}
% Retained minimal preamble for the Stacks Project chapter sites-cohomology.tex (Cohomology on Sites).
% Verbatim excerpt of the chapter preamble: Xy-pic setup (native preamble lines 14-18) and the native
% theorem declarations with the mathematical macros (native preamble lines 154-217). Copied mechanically;
% nothing added, removed or redefined. The wrapper class and the amsmath/amssymb/amsthm/graphicx base
% packages are supplied by the assembler, so the class/externaldocument part of the native preamble is omitted.
% For commutative diagrams we use Xy-pic
\usepackage[all]{xy}

% We use 2cell for 2-commutative diagrams.
\xyoption{2cell}
% Theorem environments.
%
\theoremstyle{plain}
\newtheorem{theorem}[subsection]{Theorem}
\newtheorem{proposition}[subsection]{Proposition}
\newtheorem{lemma}[subsection]{Lemma}

\theoremstyle{definition}
\newtheorem{definition}[subsection]{Definition}
\newtheorem{example}[subsection]{Example}
\newtheorem{exercise}[subsection]{Exercise}
\newtheorem{situation}[subsection]{Situation}

\theoremstyle{remark}
\newtheorem{remark}[subsection]{Remark}
\newtheorem{remarks}[subsection]{Remarks}

\numberwithin{equation}{subsection}

% Macros
%
\def\lim{\mathop{\mathrm{lim}}\nolimits}
\def\colim{\mathop{\mathrm{colim}}\nolimits}
\def\Spec{\mathop{\mathrm{Spec}}}
\def\Hom{\mathop{\mathrm{Hom}}\nolimits}
\def\Ext{\mathop{\mathrm{Ext}}\nolimits}
\def\SheafHom{\mathop{\mathcal{H}\!\mathit{om}}\nolimits}
\def\SheafExt{\mathop{\mathcal{E}\!\mathit{xt}}\nolimits}
\def\Sch{\mathit{Sch}}
\def\Mor{\mathop{\mathrm{Mor}}\nolimits}
\def\Ob{\mathop{\mathrm{Ob}}\nolimits}
\def\Sh{\mathop{\mathit{Sh}}\nolimits}
\def\NL{\mathop{N\!L}\nolimits}
\def\CH{\mathop{\mathrm{CH}}\nolimits}
\def\proetale{{pro\text{-}\acute{e}tale}}
\def\etale{{\acute{e}tale}}
\def\QCoh{\mathit{QCoh}}
\def\Ker{\mathop{\mathrm{Ker}}}
\def\Im{\mathop{\mathrm{Im}}}
\def\Coker{\mathop{\mathrm{Coker}}}
\def\Coim{\mathop{\mathrm{Coim}}}

% Boxtimes
%
\DeclareMathSymbol{\boxtimes}{\mathbin}{AMSa}{"02}

%
% Macros for moduli stacks/spaces
%
\def\QCohstack{\mathcal{QC}\!\mathit{oh}}
\def\Cohstack{\mathcal{C}\!\mathit{oh}}
\def\Spacesstack{\mathcal{S}\!\mathit{paces}}
\def\Quotfunctor{\mathrm{Quot}}
\def\Hilbfunctor{\mathrm{Hilb}}
\def\Curvesstack{\mathcal{C}\!\mathit{urves}}
\def\Polarizedstack{\mathcal{P}\!\mathit{olarized}}
\def\Complexesstack{\mathcal{C}\!\mathit{omplexes}}
% \Pic is the operator that assigns to X its picard group, usage \Pic(X)
% \Picardstack_{X/B} denotes the Picard stack of X over B
% \Picardfunctor_{X/B} denotes the Picard functor of X over B
\def\Pic{\mathop{\mathrm{Pic}}\nolimits}
\def\Picardstack{\mathcal{P}\!\mathit{ic}}
\def\Picardfunctor{\mathrm{Pic}}
\def\Deformationcategory{\mathcal{D}\!\mathit{ef}}
\usepackage{cleveref}
\crefname{originallemma}{original lemma}{original lemmas}
\Crefname{originallemma}{Original lemma}{Original lemmas}
\makeatletter
\newlabel{sites-modules-lemma-pullback-flat}{{05VD}{source}{}{sourcepointer-1}{}}
\newlabel{sites-modules-lemma-pullback-flat@cref}{{[originallemma][0][]05VD}{[1][0][]source}{}{}{}}
\makeatother
\crefname{originallemma}{original lemma}{original lemmas}
\Crefname{originallemma}{Original lemma}{Original lemmas}
\makeatletter
\newlabel{lemma-resolution-by-direct-sums-extensions-by-zero}{{077J}{source}{}{sourcepointer-2}{}}
\newlabel{lemma-resolution-by-direct-sums-extensions-by-zero@cref}{{[originallemma][0][]077J}{[1][0][]source}{}{}{}}
\makeatother
\crefname{originallemma}{original lemma}{original lemmas}
\Crefname{originallemma}{Original lemma}{Original lemmas}
\makeatletter
\newlabel{sites-modules-lemma-j-shriek-flat}{{03EV}{source}{}{sourcepointer-3}{}}
\newlabel{sites-modules-lemma-j-shriek-flat@cref}{{[originallemma][0][]03EV}{[1][0][]source}{}{}{}}
\makeatother
\crefname{originallemma}{original lemma}{original lemmas}
\Crefname{originallemma}{Original lemma}{Original lemmas}
\makeatletter
\newlabel{lemma-bounded-flat-K-flat}{{06YQ}{source}{}{sourcepointer-4}{}}
\newlabel{lemma-bounded-flat-K-flat@cref}{{[originallemma][0][]06YQ}{[1][0][]source}{}{}{}}
\makeatother
\crefname{originallemma}{original lemma}{original lemmas}
\Crefname{originallemma}{Original lemma}{Original lemmas}
\makeatletter
\newlabel{lemma-colimit-K-flat}{{06YR}{source}{}{sourcepointer-5}{}}
\newlabel{lemma-colimit-K-flat@cref}{{[originallemma][0][]06YR}{[1][0][]source}{}{}{}}
\makeatother
\crefname{originallemma}{original lemma}{original lemmas}
\Crefname{originallemma}{Original lemma}{Original lemmas}
\makeatletter
\newlabel{sites-modules-lemma-colimits-flat}{{03EU}{source}{}{sourcepointer-6}{}}
\newlabel{sites-modules-lemma-colimits-flat@cref}{{[originallemma][0][]03EU}{[1][0][]source}{}{}{}}
\makeatother
\crefname{originallemma}{original lemma}{original lemmas}
\Crefname{originallemma}{Original lemma}{Original lemmas}
\makeatletter
\newlabel{sites-modules-lemma-flat-ses}{{03EY}{source}{}{sourcepointer-7}{}}
\newlabel{sites-modules-lemma-flat-ses@cref}{{[originallemma][0][]03EY}{[1][0][]source}{}{}{}}
\makeatother
\crefname{originallemma}{original lemma}{original lemmas}
\Crefname{originallemma}{Original lemma}{Original lemmas}
\makeatletter
\newlabel{lemma-K-flat-two-out-of-three-ses}{{0G7B}{source}{}{sourcepointer-8}{}}
\newlabel{lemma-K-flat-two-out-of-three-ses@cref}{{[originallemma][0][]0G7B}{[1][0][]source}{}{}{}}
\makeatother
\crefname{originallemma}{original lemma}{original lemmas}
\Crefname{originallemma}{Original lemma}{Original lemmas}
\makeatletter
\newlabel{sites-modules-lemma-pullback-ses}{{0G6R}{source}{}{sourcepointer-9}{}}
\newlabel{sites-modules-lemma-pullback-ses@cref}{{[originallemma][0][]0G6R}{[1][0][]source}{}{}{}}
\makeatother
\crefname{originallemma}{original lemma}{original lemmas}
\Crefname{originallemma}{Original lemma}{Original lemmas}
\makeatletter
\newlabel{lemma-K-flat-flat-acyclic}{{0G7C}{source}{}{sourcepointer-10}{}}
\newlabel{lemma-K-flat-flat-acyclic@cref}{{[originallemma][0][]0G7C}{[1][0][]source}{}{}{}}
\makeatother
\title{Pullback of a K-flat complex with flat terms is again K-flat with flat terms}
\author{The Stacks Project Authors}
\date{stacks-project:sites-cohomology.tex:local-snapshot-20261011; GNU FDL 1.2 or later}
\begin{document}
\maketitle
\noindent\textit{Editorial note.} All delivered mathematical text is the author's own, copied from the pinned Stacks Project chapter file sites-cohomology.tex; no statement or proof step was rewritten, reordered or removed.\par\medskip
\noindent\textit{Editorial note.} The proof is the author's complete argument for both assertions: flatness of the pulled back terms, the reduction to truncations, and the induction using the auxiliary flatness result are all inside the counted span. The Xy-pic diagram is delivered as written.\par\medskip
\noindent\textit{Editorial note.} Every labelled reference occurring inside the counted spans points to a chapter of the Project that is not part of this card's source tree. Each is carried in assembly.external\_references as a non-load-bearing navigation pointer whose source location is the invocation inside the counted span, so no load-bearing mathematics is replaced by a stub.\par\medskip
\noindent\textit{Editorial note.} The checked chapter list for same-name collisions is the set of stacks chapters staged in this repository: chow, cohomology, divisors, etale-cohomology, flat, more-algebra, more-morphisms, morphisms, spaces-more-morphisms, varieties together with this card's chapter sites-cohomology.\par\medskip
\noindent\textit{Editorial note.} The preamble is the chapter's own retained minimal preamble (sources/stacks-sites-cohomology/retained-preamble.tex), a verbatim excerpt of the native preamble: the Xy-pic setup and the native theorem declarations with the mathematical macros.\par\medskip
\noindent\textit{Allowed standard prerequisites (editorial).}\begin{itemize}
\item K-flat complex of O-modules on a ringed topos. Same chapter, the definition of K-flatness through acyclicity of total tensor products with acyclic complexes. Supplies the property transferred by the pullback and the form in which the proof tests it.
\item Flatness of the pullback of a flat module. Modules on Sites, the result that the pullback of a flat O-module is a flat O'-module. Supplies the flatness of the terms f\textasciicircum{}*K\textasciicircum{}n, which is the first assertion of the statement.
\item Truncations of a complex and the auxiliary flatness results for them. Same chapter, the behaviour of K-flatness and flatness of terms under truncation and under the passage to a quotient, as recorded in the auxiliary lemmas of this section. Supplies the induction over the truncations with which the proof checks the acyclicity conditions defining K-flatness.
\end{itemize}\medskip
\ifdefined\hypertarget\hypertarget{sourcepointer-1}{}\fi
\noindent\textit{Source reference.} Original lemma 05VD: sources/stacks-sites-cohomology/pinned-original-source.tex, source lines 3518--3518. Navigation pointer invoked at proof line 3518; the counted span contains the pointer only and the cited result is not reproduced inside it, so the pointer is registered as a non-load-bearing navigation stub in the established form for Stacks material. No bibliography citation is involved.\par\medskip
\ifdefined\hypertarget\hypertarget{sourcepointer-2}{}\fi
\noindent\textit{Source reference.} Original lemma 077J: sources/stacks-sites-cohomology/pinned-original-source.tex, source lines 3173--3173. Navigation pointer invoked at proof line 3528; the counted span contains the pointer only and the cited result is not reproduced inside it, so the pointer is registered as a non-load-bearing navigation stub in the established form for Stacks material. No bibliography citation is involved.\par\medskip
\ifdefined\hypertarget\hypertarget{sourcepointer-3}{}\fi
\noindent\textit{Source reference.} Original lemma 03EV: sources/stacks-sites-cohomology/pinned-original-source.tex, source lines 3532--3532. Navigation pointer invoked at proof line 3532; the counted span contains the pointer only and the cited result is not reproduced inside it, so the pointer is registered as a non-load-bearing navigation stub in the established form for Stacks material. No bibliography citation is involved.\par\medskip
\ifdefined\hypertarget\hypertarget{sourcepointer-4}{}\fi
\noindent\textit{Source reference.} Original lemma 06YQ: sources/stacks-sites-cohomology/pinned-original-source.tex, source lines 3120--3120. Navigation pointer invoked at proof line 3539; the counted span contains the pointer only and the cited result is not reproduced inside it, so the pointer is registered as a non-load-bearing navigation stub in the established form for Stacks material. No bibliography citation is involved.\par\medskip
\ifdefined\hypertarget\hypertarget{sourcepointer-5}{}\fi
\noindent\textit{Source reference.} Original lemma 06YR: sources/stacks-sites-cohomology/pinned-original-source.tex, source lines 3152--3152. Navigation pointer invoked at proof line 3540; the counted span contains the pointer only and the cited result is not reproduced inside it, so the pointer is registered as a non-load-bearing navigation stub in the established form for Stacks material. No bibliography citation is involved.\par\medskip
\ifdefined\hypertarget\hypertarget{sourcepointer-6}{}\fi
\noindent\textit{Source reference.} Original lemma 03EU: sources/stacks-sites-cohomology/pinned-original-source.tex, source lines 3541--3541. Navigation pointer invoked at proof line 3541; the counted span contains the pointer only and the cited result is not reproduced inside it, so the pointer is registered as a non-load-bearing navigation stub in the established form for Stacks material. No bibliography citation is involved.\par\medskip
\ifdefined\hypertarget\hypertarget{sourcepointer-7}{}\fi
\noindent\textit{Source reference.} Original lemma 03EY: sources/stacks-sites-cohomology/pinned-original-source.tex, source lines 3542--3542. Navigation pointer invoked at proof line 3542; the counted span contains the pointer only and the cited result is not reproduced inside it, so the pointer is registered as a non-load-bearing navigation stub in the established form for Stacks material. No bibliography citation is involved.\par\medskip
\ifdefined\hypertarget\hypertarget{sourcepointer-8}{}\fi
\noindent\textit{Source reference.} Original lemma 0G7B: sources/stacks-sites-cohomology/pinned-original-source.tex, source lines 3097--3097. Navigation pointer invoked at proof line 3544; the counted span contains the pointer only and the cited result is not reproduced inside it, so the pointer is registered as a non-load-bearing navigation stub in the established form for Stacks material. No bibliography citation is involved.\par\medskip
\ifdefined\hypertarget\hypertarget{sourcepointer-9}{}\fi
\noindent\textit{Source reference.} Original lemma 0G6R: sources/stacks-sites-cohomology/pinned-original-source.tex, source lines 3558--3558. Navigation pointer invoked at proof line 3558; the counted span contains the pointer only and the cited result is not reproduced inside it, so the pointer is registered as a non-load-bearing navigation stub in the established form for Stacks material. No bibliography citation is involved.\par\medskip
\ifdefined\hypertarget\hypertarget{sourcepointer-10}{}\fi
\noindent\textit{Source reference.} Original lemma 0G7C: sources/stacks-sites-cohomology/pinned-original-source.tex, source lines 3385--3385. Navigation pointer invoked at proof line 3565; the counted span contains the pointer only and the cited result is not reproduced inside it, so the pointer is registered as a non-load-bearing navigation stub in the established form for Stacks material. No bibliography citation is involved.\par\medskip
\begin{lemma}
\label{lemma-pullback-K-flat}
Let $f : (\Sh(\mathcal{C}'), \mathcal{O}') \to (\Sh(\mathcal{C}), \mathcal{O})$
be a morphism of ringed topoi. Let $\mathcal{K}^\bullet$ be a K-flat complex
of $\mathcal{O}$-modules whose terms are flat $\mathcal{O}$-modules. Then
$f^*\mathcal{K}^\bullet$ is a K-flat complex of $\mathcal{O}'$-modules whose
terms are flat $\mathcal{O}'$-modules.
\end{lemma}

\begin{proof}
The terms $f^*\mathcal{K}^n$ are flat $\mathcal{O}'$-modules by
Modules on Sites, Lemma \ref{sites-modules-lemma-pullback-flat}.
Choose a diagram
$$
\xymatrix{
\mathcal{K}_1^\bullet \ar[d] \ar[r] &
\mathcal{K}_2^\bullet \ar[d] \ar[r] & \ldots \\
\tau_{\leq 1}\mathcal{K}^\bullet \ar[r] &
\tau_{\leq 2}\mathcal{K}^\bullet \ar[r] & \ldots
}
$$
as in Lemma \ref{lemma-resolution-by-direct-sums-extensions-by-zero}.
We will use all of the properties stated in the
lemma without further mention. Each $\mathcal{K}_n^\bullet$ is a bounded
above complex of flat modules, see
Modules on Sites, Lemma \ref{sites-modules-lemma-j-shriek-flat}.
Consider the short exact sequence of complexes
$$
0 \to \mathcal{M}^\bullet \to
\colim \mathcal{K}_n^\bullet \to
\mathcal{K}^\bullet \to 0
$$
defining $\mathcal{M}^\bullet$. By Lemmas \ref{lemma-bounded-flat-K-flat} and
\ref{lemma-colimit-K-flat} the complex $\colim \mathcal{K}_n^\bullet$
is K-flat and by Modules on Sites, Lemma \ref{sites-modules-lemma-colimits-flat}
it has flat terms. By Modules on Sites, Lemma \ref{sites-modules-lemma-flat-ses}
$\mathcal{M}^\bullet$ has flat terms, by
Lemma \ref{lemma-K-flat-two-out-of-three-ses}
$\mathcal{M}^\bullet$ is K-flat, and by the long exact
cohomology sequence $\mathcal{M}^\bullet$ is acyclic (because
the second arrow is a quasi-isomorphism). The pullback
$f^*(\colim \mathcal{K}_n^\bullet) = \colim f^*\mathcal{K}_n^\bullet$
is a colimit of bounded below complexes of flat $\mathcal{O}'$-modules
and hence is K-flat (by the same lemmas as above).
The pullback of our short exact sequence
$$
0 \to f^*\mathcal{M}^\bullet \to
f^*(\colim \mathcal{K}_n^\bullet) \to
f^*\mathcal{K}^\bullet \to 0
$$
is a short exact sequence of complexes by
Modules on Sites, Lemma \ref{sites-modules-lemma-pullback-ses}.
Hence by Lemma \ref{lemma-K-flat-two-out-of-three-ses}
it suffices to show that $f^*\mathcal{M}^\bullet$
is K-flat. This reduces us to the case discussed in the next paragraph.

\medskip\noindent
Assume $\mathcal{K}^\bullet$ is acyclic as well as K-flat and
with flat terms. Then Lemma \ref{lemma-K-flat-flat-acyclic}
guarantees that all terms of $\tau_{\leq n}\mathcal{K}^\bullet$
are flat $\mathcal{O}$-modules. We choose a diagram as above and
we will use all the properties proven above for this diagram.
Denote $\mathcal{M}_n^\bullet$ the kernel of the map of complexes
$\mathcal{K}_n^\bullet \to \tau_{\leq n}\mathcal{K}^\bullet$
so that we have short exact sequences of complexes
$$
0 \to \mathcal{M}_n^\bullet \to \mathcal{K}_n^\bullet \to
\tau_{\leq n}\mathcal{K}^\bullet \to 0
$$
By Modules on Sites, Lemma \ref{sites-modules-lemma-flat-ses}
we see that the terms of the complex $\mathcal{M}_n^\bullet$ are flat.
Hence we see that $\mathcal{M} = \colim \mathcal{M}_n^\bullet$
is a filtered colimit of bounded below complexes of flat modules
in this case. Thus $f^*\mathcal{M}^\bullet$ is K-flat
(same argument as above) and we win.
\end{proof}

\end{document}
